% Generated draft: requires mathematical and source review.
\documentclass[11pt]{article}
\usepackage[margin=25mm]{geometry}
\usepackage{fontspec}
\setmainfont{Latin Modern Roman}
\usepackage{amsmath,amssymb,mathtools,unicode-math}
\setmathfont{Latin Modern Math}
\usepackage{enumitem,needspace,xurl,hyperref}
\setlength{\parindent}{0pt}
\setlength{\parskip}{5pt}
\setlength{\emergencystretch}{3em}
\newcommand{\N}{\mathbb{N}}
\newcommand{\Z}{\mathbb{Z}}
\newcommand{\Q}{\mathbb{Q}}
\newcommand{\R}{\mathbb{R}}
\newcommand{\C}{\mathbb{C}}
\newcommand{\F}{\mathbb{F}}
\newcommand{\A}{\mathbb{A}}
\newcommand{\PP}{\mathbb{P}}
\newcommand{\E}{\mathbb{E}}
\newcommand{\eps}{\varepsilon}
\newcommand{\dd}{\,\mathrm{d}}
\begin{document}
\section*{UnsolvedMath ID: 30001860}
\textbf{Source problem code:} catalogue entry
\subsection{Definitions and mathematical statement}
Let \(G=\operatorname{Class}(n,q)\) denote a finite classical group with natural module \(V\) of dimension \(n\) over a field of odd order \(q\), as stipulated in the supplied record. The notation \(\operatorname{Class}(n,q)\) does not specify the precise classical-group family or version of the group; the definitions below require a specified linear action on \(V\). Dimensions are taken over the defining field of that module.

For \(g\in G\), let \(|g|\) denote its order. If \(|g|\) is even, define the involution
\[
 t_g=g^{|g|/2},\qquad \operatorname{Fix}_V(t_g)=\{v\in V:t_gv=v\}.
\]
Define
\[
 Q=\left\{g\in G: |g|\text{ is even and }\frac n3\leq \dim\operatorname{Fix}_V(t_g)<\frac{2n}{3}\right\},
 \qquad p(G)=\frac{|Q|}{|G|}.
\]
The lower endpoint is included and the upper endpoint is excluded. Since the dimension is an integer, the dimension condition is equivalently
\[
 \left\lceil\frac n3\right\rceil\leq\dim\operatorname{Fix}_V(t_g)\leq\left\lceil\frac{2n}{3}\right\rceil-1.
\]

The supplied cleaned question asks for the asymptotic behavior of \(p(G)\), and in particular whether it has order \(1/\log n\). One possible precise interpretation, after fixing a classical-group family and an odd field order \(q\), is whether there exist constants \(c,C>0\) and \(n_0\) such that
\[
 \frac{c}{\log n}\leq p\bigl(\operatorname{Class}(n,q)\bigr)\leq\frac{C}{\log n}
\]
for every admissible dimension \(n\geq n_0\). This is an explicitly conditional interpretation: the supplied record does not specify whether \(q\) is fixed or varies with \(n\), which classical families are covered, or whether the constants must be uniform in \(q\) and the family.

There is also an unresolved discrepancy in the target. The supplied original formulation asks instead whether
\[
 \frac{|Q|}{|G|}=O(\log n).
\]
The supplied verification note says that the printed upper-bound notation appears to omit a reciprocal and that the cleaned formulation follows the order suggested by an adjacent lower bound. Neither the relevant passage nor that lower bound is reproduced in the record. Thus the reciprocal correction and the interpretation as a two-sided asymptotic order remain source qualifications requiring review; the two formulations must not be identified without qualification.
\subsection{Short English statement}
For a finite classical group acting on its natural module, consider the even-order elements whose associated involution fixes at least one third but strictly less than two thirds of the module's dimension. Determine their proportion as the dimension grows. The cleaned record suggests order \(1/\log n\), but the original bound, the precise group family, and the required uniformity are unresolved.
\subsection{Sources}
\begin{itemize}
\item Local UnsolvedMath registry, supplied complete source record: permanent ID \(30001860\), source problem code \texttt{catalogue entry}, title ``Asymptotic Proportion of \(Q\) in Classical Groups,'' in the Oberwolfach Reports Open Problems collection. No individual proposer is supplied; the proposed year is recorded as 2011.
\item Supplied report citation: \emph{Computational Group Theory (2012)}, Oberwolfach Reports, DOI \texttt{10.4171/owr/2011/37}, pp.\ 2113--2161. Supplied source URL and literature-source link labeled ``OWR report 37'': \url{https://doi.org/10.4171/owr/2011/37}. The citation's 2012 label and the record's 2011 year are retained as supplied; their relationship has not been independently checked.
\item The supplied record labels the statement \texttt{corrected\_verified} and reports that the official report was checked to restore the classical-group setup and asymptotic question. It also qualifies the reciprocal correction as an interpretation of apparently omitted notation. These are supplied verification claims, not independent verification performed for this draft.
\item Supplied dated literature curation, checked at \texttt{2026-08-21}: status ``open''; the assessment states that no authoritative resolution of the stated asymptotic bound was found and that the report poses it as an open problem. The listed source checked is ``OWR report 37,'' \url{https://doi.org/10.4171/owr/2011/37}. This is supplied dated curation, not an independently verified literature assessment; time-sensitive status and cited claims require independent review.
\end{itemize}
\end{document}
